\documentclass[../main]{subfiles}
\begin{document}
\ifSubfilesClassLoaded{\setcounter{page}{4}}{}
\section{Inventiamo gli «occhiali-bicicletta»}
\label{sec:01_inventiamo_occhiali}

Cos'è il comunismo? Quali sono le sue origini, dove ci esorta a dirigere i nostri passi, e chi, precisamente, ci invita a seguirlo? Queste domande, apparentemente semplici, richiedono risposte puntuali, come la Storia ci ha ampiamente dimostrato, non una tantum, bensì ogni qualvolta si ripresentino. Come affermò il drammaturgo Grigorij Gorin, attraverso le parole del «celebre Munchausen», «l'amore è un teorema che deve essere dimostrato ogni giorno». E per noi il comunismo rappresenta una sfida analoga, sebbene si tratti di un teorema (o, più precisamente, di una teoria) di portata più ampia e di maggiore complessità. Dobbiamo quindi imparare, innanzitutto, a elencare con precisione, nel punto dedicato ai «dati» di questo teorema, tutti gli elementi che effettivamente lo compongono, distinguendoli da quelli che appartengono invece al regno dell'immaginazione.

Le critiche rivolte al termine stesso “comunismo” riflettono il declino storico che ha colpito la nostra società. La sua svalutazione, e persino la sua trasformazione in un’espressione denigratoria, fenomeni che si sono manifestati dopo il collasso dell’Unione Sovietica, sono il risultato dell’ascesa di forze sociali apertamente avverse al comunismo.

Ben prima del collasso dell'Unione Sovietica, il concetto aveva già iniziato a perdere chiarezza, tanto che alcuni teorici auspicavano una distinzione all'interno del sistema economico e ideologico sovietico, tracciando una linea di demarcazione per evitare di confondere l'inerzia capitalistica con una presunta tendenza comunista. Successivamente, poiché la società venne improvvisamente catapultata in una restaurazione del capitalismo in un contesto in cui la tradizione del pensiero comunista era quasi del tutto svanita, e la questione si ridusse alla sua mera sopravvivenza, i problemi del comunismo e del movimento che ad esso aspirava cessarono di essere oggetto di studio specifico. Tutti gli interessi si concentrarono sul capitalismo in piena espansione.Sia i suoi fautori che i suoi detrattori lo consideravano una figura tangibile, mentre il comunismo e il movimento ad esso legato apparivano ormai relegati al passato, oppure, nel caso dei suoi sostenitori, relegati a un futuro così remoto da sembrare irraggiungibile, rendendo quindi prioritario focalizzarsi sull'esistenza presente. Il comunismo, inteso come questione teorica, ha smesso di essere rilevante da decenni.

Nella cultura di massa, estranea alla speculazione teorica, la coscienza della classe dominante è stata (e continua a essere) attivamente riprodotta. Questa classe, pervasa dal timore e dall’odio verso il comunismo, trasforma questi sentimenti in molteplici forme di miti sociali, sui quali, nelle società post-sovietiche, sta crescendo una seconda generazione. Tuttavia, tali rappresentazioni mitizzate del comunismo non hanno alcuna attinenza con il suo concetto originario. Qui si cela il primo e più rilevante punto di “contrasto” (che, in realtà, non lo è). Gli abitanti dei paesi che hanno abbracciato il capitalismo, non in seguito a un’imposizione esterna, ma attraverso un processo interno, semplicemente non conoscono più l’essenza del vero comunismo, ossia il comunismo scientifico, l’anima stessa dell’ideale comunista. E, nonostante ciò, la questione del comunismo si ripropone come un problema concreto e attuale.Anche quella teoria organica, la cui messa al bando e il cui discredito hanno contribuito in maniera determinante al ritorno al capitalismo nell’URSS e negli altri paesi del blocco socialista, necessita di essere recuperata e presentata in forma chiara e accessibile, affinché possa raggiungere coloro i cui interessi essa effettivamente tutela. Non si tratta unicamente degli operai che mostrano indifferenza politica, ma anche di quanti si definiscono comunisti o, perlomeno, di sinistra, ovvero di coloro che si riconoscono, a livello politico, come sostenitori del comunismo. Questi ultimi tendono a concepire il comunismo come una sorta di sistema sociale “positivo”, interpretando tale “positività” soprattutto a livello di percezioni sensoriali, spesso per contrasto, come l’assenza delle miserie e degli orrori che affliggono la società capitalista contemporanea. Tuttavia, questi sostenitori non sono in grado di fornire una risposta compiuta alla domanda “che cos’è il comunismo?”, nemmeno a se stessi.

Le circostanze attuali ci inducono a dare vita a questa pubblicazione, concepita come una sorta di "strumento di analisi", uno strumento che ci consenta di esaminare il comunismo con occhi nuovi e di coglierne l'essenza più profonda. Nei confronti di coloro che si oppongono al comunismo, è opportuno, a nostro avviso, richiedere che formulino obiezioni pertinenti e razionali, anziché semplici e prive di fondamento, di fronte a questa dottrina, ora intesa in una prospettiva scientifica.

Nell’anno in cui si celebrava il centenario della Rivoluzione d’Ottobre, le interpretazioni sul ruolo e sul comportamento dei bolscevichi nel 1917, emerse nel dibattito pubblico, assumevano sfumature così diverse da dar luogo a una polarizzazione inevitabile: da un lato, l’immagine di eroi nobili che “salvarono lo Stato”; dall’altro, quella di banditi animati dall’odio e dalla volontà di distruggere lo stesso (la statualità). In tale contesto di estremi contrapposti, tuttavia, un aspetto cruciale è stato trascurato: i bolscevichi, e in particolare Lenin, erano degli intellettuali, e le loro azioni furono guidate non solo da desideri e preferenze, ma anche da una solida teoria, frutto dell’intero percorso evolutivo del pensiero umano. I detrattori del socialismo e del comunismo tendono a definire la Rivoluzione d’Ottobre e la successiva costruzione socialista come un esperimento condotto dai bolscevichi sulle popolazioni… Ebbene, tali oppositori si avvicinano, forse, più della loro controparte, alla comprensione della reale natura di quell’evento, che fu profondamente radicato in una specifica elaborazione teorica.

La Grande Rivoluzione d’Ottobre fu, in fondo, un esperimento, a meno che non si intenda il termine “esperimento” in una accezione banale, come una mera manovra superficiale. Data la sua natura inedita, essa non poteva essere altro che un tentativo, condotto secondo le regole di ogni rigorosa indagine scientifica: infatti, non esiste esperimento privo di una precisa ipotesi. E l’ipotesi, formulata da Lenin e dai suoi seguaci, e fondata sull’eredità di Marx ed Engels, era quella di una Rivoluzione Mondiale. Iniziata in Russia, la rivoluzione proletaria avrebbe dovuto estendersi ai paesi europei industrializzati, dando inizio alla loro trasformazione in una Repubblica Mondiale dei Soviet.È vero che non tutte le ipotesi scientifiche sopravvivono alla verifica sperimentale nella loro formulazione originaria (sebbene l’idea fondamentale della validità delle rivoluzioni proletarie mantenga, in linea generale, la sua attualità anche ai giorni nostri). Ciononostante, anche questo tipo di esperimenti – condotti seguendo le tracce di un’ipotesi che si è rivelata solo parzialmente corretta – contribuiscono al progresso della scienza (e, in questo caso, del socialismo).

La metafora che ci viene proposta apre la strada a una diversa prospettiva sulla rivoluzione, sul comunismo di guerra, sulla NEP, sull’industrializzazione e via dicendo, esortandoci a considerare la società contemporanea e le direttrici del suo sviluppo attraverso la lente del comunismo, inteso come scienza fondamentale all’interno del sistema delle scienze sociali. Solo in questo modo potremo affrontare con razionalità le sfide del mondo moderno. Affinché la classe rivoluzionaria possa liberarsi dalla sua condizione di subalternità allo sfruttamento, superare la propria identità di classe oppressa nel sistema capitalistico e assumere una posizione autonoma, è necessario che essa comprenda e modifichi in modo consapevole i rapporti sociali, orientandoli verso l’eliminazione di ogni forma di oppressione e sfruttamento, verso un futuro comunista.

È per questo che le teorie comuniste, oggetto di sardonica derisione da parte di Majakovskij – che le definì con l’immagine degli «occhiali-bicicletta» – potranno rivelarsi utili solo a condizione che riscopriamo la pazienza necessaria per «reinventarle», ovvero per riprodurne i principi fondamentali. Dobbiamo ripercorrere il cammino intrapreso da Marx, dalla Comune all’Internazionale, e seguire l’evoluzione del pensiero leniniano fino alla Rivoluzione d’Ottobre, per poi analizzare l’intera esperienza sovietica, senza alterare i «mezzi di trasporto», vale a dire la teoria rivoluzionaria.

Rinnovare significa dunque appropriarci dell’eredità della scienza comunista, adattandola alle condizioni storiche attuali, quelle che ci sono state trasmesse.
\end{document}
